% =============================================================
% §1 Introduction
% =============================================================
\section{Introduction}
\label{sec:intro}

Generative AI workloads---also termed AI-generated content
(AIGC) workloads---dominated by large language model (LLM)
inference, have become the fastest-growing class of compute in
modern data centers: their energy cost has grown by orders of
magnitude over five years~\cite{patterson2021carbon}, and
inference now consumes more aggregate FLOPs than
training~\cite{patel2024splitwise}. As LLM services move into
latency-sensitive applications, the systems community is
rethinking how inference workloads should be \emph{scheduled}
across heterogeneous compute pools.

This paper studies \textbf{cloud-edge collaborative LLM inference
scheduling}, a deployment paradigm gaining traction for three
reasons: (i)~\textbf{cost}, by serving common-case queries on
cheaper edge accelerators while reserving cloud GPUs for tail
requests; (ii)~\textbf{privacy}, by keeping prompt data on-premises
whenever possible; and (iii)~\textbf{tail latency}, by amortizing
first-token latency through edge proximity.
\emph{Collaborative} here denotes deployment-level
cooperation---cloud and edge tiers jointly serving one
workload---not multi-agent decision-making: every scheduler we
study is a single centralized policy. Our objectives
(\S\ref{sec:sched-problem}) directly capture tail latency, via
service-level-objective (\slo) attainment, and cost, via energy
per token---the dominant marginal serving cost; privacy
motivates \emph{deploying} cloud-edge pools but is orthogonal
to per-request placement (it can be encoded as feasibility
masks, \S\ref{sec:mdp}, which we leave to future
work). The scheduling
problem---deciding which server should serve each request---now
becomes a critical performance lever, distinct from intra-server
batching engines such as
vLLM~\cite{kwon2023vllm}, Orca~\cite{yu2022orca}, and
Sarathi-Serve~\cite{agrawal2024sarathi}.

\subsection{What makes LLM inference scheduling different}

Existing cluster schedulers---list-scheduling heuristics such
as Heterogeneous Earliest Finish Time (HEFT), metaheuristics
such as genetic algorithms (GA) and particle swarm optimization
(PSO), and learning-based variants like
Decima~\cite{mao2019decima} and
RLScheduler~\cite{zhang2020rlscheduler}---were designed for
generic directed-acyclic-graph (DAG) workloads such as Spark
jobs or high-performance-computing (HPC) batch tasks. They do not
model the five physical characteristics that fundamentally
distinguish LLM inference:

\begin{enumerate}
\item \textbf{Model weight residency.} LLM weights (14--70~GB
  for LLaMA-7B/13B/70B~\cite{touvron2023llama}) reside
  persistently in GPU memory; switching models costs 5--60~s
  of cold load~\cite{fu2024serverlessllm}---orders of
  magnitude beyond any per-request computation.

\item \textbf{KV-cache locality.} Each request's key-value (KV)
  cache is bound to the GPU where prefill ran; migrating it
  moves hundreds of MB per request, often dominating decode
  latency~\cite{patel2024splitwise,kwon2023vllm}.

\item \textbf{Continuous batching.} Iteration-level batching
  serves $N$ same-model requests in
  ${\approx}(1 + (N-1)\alpha)\,T_{\text{solo}}$ time---a
  $5$--$8\times$ throughput
  gain~\cite{yu2022orca,kwon2023vllm} that materializes only
  if the scheduler co-locates same-model requests.

\item \textbf{Memory-bandwidth-bound decode.} Unlike
  compute-bound prefill, decode is bound by
  high-bandwidth-memory (HBM) bandwidth: ${\approx}20$~ms per
  token on A100-class
  GPUs~\cite{patel2024splitwise,agrawal2024sarathi} regardless
  of compute capacity.

\item \textbf{Two-phase request lifecycle.} Each request
  decomposes into Prefill$\,\to\,$Decode with strict dependency
  and shared KV state~\cite{zhong2024distserve}, violating the
  independent-stage assumption of most DAG schedulers.
\end{enumerate}

These five physics create trade-offs absent from generic
workloads: consolidating same-model requests (for batching)
versus spreading them (against contention); dispatching decode
to a fast-but-distant server versus its sibling prefill server
(no KV migration). \textbf{A scheduler that ignores AIGC
physics cannot exploit these trade-offs}.

\subsection{Why na\"ive approaches fail}

A natural hypothesis is that an off-the-shelf
reinforcement-learning (RL) scheduler would \emph{implicitly}
learn AIGC-aware behavior from base features (compute, queue
length, memory). We set out to test this hypothesis; the
empirical picture is more subtle (Fig.~\ref{fig:pareto}):

\begin{itemize}
\item Generic load-balancing baselines (LeastLoaded,
  ShortestQueue), optimal-DAG heuristics (HEFT, GA, PSO),
  a PerLLM-style constrained-bandit
  scheduler~\cite{yang2024perllm}, and
  even modern deep-reinforcement-learning (DRL) baselines
  (A3C-R2N2~\cite{tuli2022a3cr2n} and
  a Decima-inspired graph-neural-network (GNN) variant we
  implemented~\cite{mao2019decima}) all cluster in the same
  Pareto region: \slo{} attainment
  around 19--24\% and energy
  efficiency around 2.8--3.0 J/token.

\item Our AIGC-aware proximal-policy-optimization (PPO)
  scheduler, sharing the actor-critic
  family and equivalent pretraining budget with A3C-R2N2 but
  with explicit AIGC state features and reward shaping, achieves
  a \textbf{distinctly better Pareto point}: \slo{} 30.2\%
  ($+28\%$ relative) and energy 2.61~J/token ($-7\%$ relative),
  with $p<0.05$ on \slo{} against all five strongest
  single-objective baselines. The strongest baseline
  overall---a knee-point NSGA-II search we also
  implement---approaches our \slo{} at the canonical point
  but is beaten significantly at moderate load ($p<0.01$),
  pays more energy in every configuration, and spends
  $5\times$ more compute at dispatch time.

\item However, our 12-component ablation
  (Table~\ref{tab:ablation}, Fig.~\ref{fig:ablation}) reveals a
  paradox: removing any single AIGC component (warm reward,
  batch reward, affinity reward, AIGC state features, even all
  of them jointly) leaves performance essentially unchanged.
  Only continuous batching simulation and adequate pretraining
  stand out as significant components.
\end{itemize}

\subsection{Reconciliation and contributions}

These observations are reconciled by recognizing that
\textbf{AIGC-aware scheduling is a holistic system contribution,
not a clever reward component}. Once a scheduler has
(i)~continuous batching physics in its environment,
(ii)~adequate pretraining, and (iii)~some AIGC
awareness in either state or reward, the PPO backbone extracts
equivalent signal from any of the AIGC priors---they are
mutually substitutable. The Pareto improvement is
real, but it comes from the \emph{combination}, not from
individual reward tricks.

\smallskip
\noindent This paper makes four contributions:

\noindent\textbf{(C1) An AIGC inference simulation platform}
modeling all five physics: weight residency with
least-recently-used (LRU) eviction, two-phase prefill/decode
tasks with KV-cache accounting, admission-time continuous
batching, and memory-bandwidth floors calibrated to vLLM
measurements. Open-sourced with reproducibility manifests.

\noindent\textbf{(C2) An AIGC-aware PPO scheduler} (RL) that
exposes AIGC physics to the policy via state features
(loaded-model indicator, batch occupancy, sibling-server
affinity, KV-cache size) and reward shaping (warm, batch,
affinity bonuses; cloud-overuse penalty), achieving
\textbf{Pareto dominance on mean performance} over
ten strong baselines
in \slo{} attainment and energy efficiency.

\noindent\textbf{(C3) A systematic 12-component ablation study}
revealing which AIGC abstractions actually contribute. The
striking null result---no single AIGC reward or state
component is individually significant---cautions against
attributing gains to specific reward designs without joint
ablation.

\noindent\textbf{(C4) An empirical ``regime'' characterization}
showing that AIGC-aware RL achieves strict Pareto dominance only
in the \textbf{moderate-load operating regime}
($\lambda \in [1, 2]$ req/s with uniform model mix), while
maintaining the lowest energy per token across
all 30 tested configurations, with margins over the strongest
baseline growing with load---a nuanced finding that contradicts
the ``always wins'' claims common in DRL scheduling literature.

\smallskip
The remainder of the paper covers background
(\S\ref{sec:background}), the formal problem
(\S\ref{sec:formulation}), simulator and scheduler design
(\S\ref{sec:system}), evaluation (\S\ref{sec:eval}),
discussion (\S\ref{sec:discuss}), related work
(\S\ref{sec:related}), and conclusions
(\S\ref{sec:conclusion}).
